\section{因果贝叶斯网络}

把有向无环图解释为独立假设的载体，并不必然意味着因果；事实上，只要存在沿某一变量次序的递归独立关系，这种解释便成立，而该次序未必是因果的或时间的。然而，DAG 模型在统计学与人工智能应用中之所以无处不在，主要还是（而且往往是不自觉地）源于其因果解释——即把它视为一个由若干过程组成的系统，每个家族对应一个过程，足以说明观测数据的生成。正是这种因果解释，说明了为什么 DAG 模型极少按不尊重时间与因果方向的变量次序来使用。

围绕因果信息而非关联信息来建立 DAG 模型，有几方面的好处。第一，构造模型时所要作的判断更有意义、更易获取，因而更可靠。读者若试着按次序 $(X_5,X_1,X_3,X_2,X_4)$ 为图 1.2 中的关联构造一个 DAG 表示，就能体会这一点。这样的练习不仅说明有些独立关系比另一些更容易被心智把握，而且说明条件独立判断只有锚定在更基本的建筑块——例如因果关系——之上时才可获取（因而可靠）。在图 1.2 的例子中，我们之所以愿意断言一旦知道 $X_4$（即路面是否湿）后 $X_5$ 就与 $X_2$、$X_3$ 独立，是因为我们可以轻易地把这一断言转译为涉及因果关系的说法：雨和洒水器对路面打滑的影响是以路面湿度为中介的。没有因果连线支撑的依赖关系会被视为古怪或虚假，甚至被冠以“悖论”之名（参见 \S 1.2.3 关于伯克森悖论的讨论）。

本书将有多次机会展示因果知识相对于关联知识的优先地位。在极端情形下，我们会看到人们倾向于完全忽略概率信息而专注于因果信息（见 \S 6.1.4）。\footnote{Tversky and Kahneman (1980) 关于概率判断中因果偏误的实验，是支持这一观察的另一组证据。例如，多数人相信“已知母亲蓝眼，则女儿蓝眼”的概率大于反过来；而这两个概率实际上相等。}这就对统计学中图模型的主流范式（Wermuth and Lauritzen 1990; Cox and Wermuth 1996）提出了质疑：该范式认为条件独立假设是表达实质性知识的主要载体。\footnote{作者本人也曾像统计学界的同行一样，主张条件独立的中心地位；见 Pearl (1988b, p.~79)。}看来，既然条件独立判断是所储存的因果关系的副产品，那么直接挖掘并表示这些关系，将是表达我们所知或所信之世界的一种更自然、更可靠的方式。这正是因果贝叶斯网络背后的哲学。

以因果关系为基础构造贝叶斯网络的第二个好处——它对理解因果组织至关重要——是有能力表示并响应外部变化或自发变化。环境中机制的任何局部重构，都只需作微小修改，便可翻译为网络拓扑的同构重构。例如，要表示图 1.2 故事中洒水器损坏，我们只需删去网络中所有与节点“洒水器”相连的连线。要表示“下雨就关掉洒水器”的策略，只需在“雨”与“洒水器”之间加一条连线，并修改 $P(x_3\mid x_1,x_2)$。若网络不是沿因果方向、而是沿（比如说）次序 $(X_5,X_1,X_3,X_2,X_4)$ 构造的，这类改动就需要大得多的重塑工作。这种重塑的灵活性，或许正是划分“审议型”与“反应型”智能体的要素，使前者能够即时应对新情境，而无需训练或适应。

\subsection{因果网络作为干预的神谕}

这种灵活性的根源在于一个假设：网络中每一对父子关系都代表一个稳定而自主的物理机制——换言之，改变其中一个关系而不改变其他关系是可以设想的。以这种模块化配置组织知识，使人能够以最少的额外信息预测外部干预的效果。的确，因果模型（假设其有效）比概率模型的信息量大得多。联合分布告诉我们事件有多可能，以及概率将如何随后的观测而变化；而因果模型还告诉我们，这些概率将如何因外部干预而变化——例如政策分析、治疗管理或日常活动规划中所遇到的那些干预。这类变化无法从联合分布中推出，即便该分布已被完全指定。

模块性与干预之间的联系如下。我们不必为众多可能干预中的每一个分别指定新的概率函数，而只需指定干预所蕴含的直接变化；凭借自主性，我们假定该变化是局部的，不会蔓延到所指定机制之外的机制。一旦知道了被干预所改变的机制是哪一个，以及改变的性质，就可以通过修改 \textup{(1.33)} 中相应的因子、并用修改后的乘积计算新的概率函数，来预测干预的总体效果。例如，要在图 1.2 的网络中表示“打开洒水器”这一行动，我们删去连线 $X_1\to X_3$，并令 $X_3$ 取值“开”。该操作所得的图见图 1.4，其余变量上所得的联合分布为
\begin{equation}
P_{X_3=\mathrm{On}}(x_1,x_2,x_4,x_5) = P(x_1)\,P(x_2\mid x_1)\,P(x_4\mid x_2,X_3=\mathrm{On})\,P(x_5\mid x_4), \tag{1.36}
\end{equation}
其中右端各因子凭借自主性均与 \textup{(1.34)} 中相同。

删去因子 $P(x_3\mid x_1)$，表达的正是这样的理解：无论行动之前季节与洒水器之间存在何种关系，在我们实施行动期间该关系都不再起作用。一旦我们实际打开洒水器并使之保持开启，决定洒水器状态的便是一个新机制（季节在其中没有发言权）。

注意行动 $\dox{X_3=\mathrm{On}}$ 与观测 $X_3=\mathrm{On}$ 之间的区别。后者的效果由普通的贝叶斯条件化得到，即 $P(x_1,x_2,x_4,x_5\mid X_3=\mathrm{On})$；而前者的效果则是对一个删去了连线 $X_1\to X_3$ 的切除图取条件。这确实映射了“看见”与“行动”之别：在观测到洒水器开着之后，我们希望推断季节是干燥的、大概没有下雨，等等；而在评估“打开洒水器”这一设想行动的效果时，不应作出任何此类推断。

因果网络预测行动效果的能力，当然要求在构造网络时采用更强的一组假设；这些假设立足于因果知识（而不只是关联知识），并确保系统会按照自主性原则对干预作出响应。这些假设凝聚在下面关于因果贝叶斯网络的定义之中。

\begin{figure}[htbp]
\centering
\begin{tikzpicture}[causal,scale=1.0,every node/.style={transform shape}]
  \node[vtx] (X1) at (0,0) {$X_1$};  \node[right=1.5mm of X1,font=\small] {SEASON};
  \node[vtxb] (X3) at (-2.0,-1.35) {$X_3$};
  \node[left=1.5mm of X3,font=\small,align=right] {SPRINKLER\\= ON};
  \node[vtx] (X2) at (2.0,-1.35) {$X_2$};  \node[right=1.5mm of X2,font=\small] {RAIN};
  \node[vtx] (X4) at (0,-2.7) {$X_4$};  \node[right=1.5mm of X4,font=\small] {WET};
  \node[vtx] (X5) at (0,-4.05) {$X_5$}; \node[right=1.5mm of X5,font=\small] {SLIPPERY};
  \draw[ed] (X1) -- (X2);
  \draw[ed] (X3) -- (X4);
  \draw[ed] (X2) -- (X4);
  \draw[ed] (X4) -- (X5);
\end{tikzpicture}
\caption{“打开洒水器”这一行动的网络表示。}
\label{fig:1.4}
\end{figure}

\begin{definition}[因果贝叶斯网络]
设 $P(v)$ 是变量集 $V$ 上的概率分布，$P_x(v)$ 表示把变量子集 $X$ 设为常数 $x$ 的干预 $\dox{X=x}$ 所导致的分布。\footnote{记号 $P_x(v)$ 在随后各章将改用 $P(y\mid\doxp(x))$ 与 $P(y\mid\hat{x})$，以便于代数运算。}用 $P^*$ 表示全体干预分布 $P_x(v),\ X\subseteq V$ 的集合，其中包括 $P(v)$，它表示无干预（即 $X=\varnothing$）。DAG $G$ 称为与 $P^*$ 相容的\emph{因果贝叶斯网络}，当且仅当对每个 $P_x\in P^*$ 下列三个条件都成立：
\begin{enumerate}[label=(\roman*)]
\item $P_x(v)$ 相对于 $G$ 是马尔可夫的；
\item 只要 $v_i$ 与 $X=x$ 相容，就有 $P_x(v_i)=1$；
\item 只要 $\pa_i$ 与 $X=x$ 相容，就有 $P_x(v_i\mid\pa_i)=P(v_i\mid\pa_i)$，即每个 $P(v_i\mid\pa_i)$ 对不涉及 $V_i$ 的干预保持不变。
\end{enumerate}
\end{definition}

定义 1.3.1 对干预空间 $P^*$ 施加了约束，使我们能够以单个贝叶斯网络 $G$ 的形式经济地编码这个庞大的空间。这些约束使我们能够把任一干预 $\dox{X=x}$ 所导致的分布 $P_x(v)$ 计算为一个截断分解
\begin{equation}
P_x(v) = \prod_{\{i\,:\,V_i\notin X\}} P(v_i\mid\pa_i)\qquad\text{对一切与 } x \text{ 相容的 } v, \tag{1.37}
\end{equation}
它由条件 (i)—(iii) 推出（并且蕴含这些条件），从而为 $G$ 上的家族删除程序——如 \textup{(1.36)} 所示——提供了依据。不难证明，只要 $G$ 关于 $P^*$ 是因果贝叶斯网络，下列两条性质必定成立。

\textbf{性质 1}\quad 对一切 $i$，
\begin{equation}
P(v_i\mid\pa_i) = P_{\pa_i}(v_i). \tag{1.38}
\end{equation}

\textbf{性质 2}\quad 对一切 $i$，以及对与 $\{V_i,\PA_i\}$ 不相交的任一变量子集 $S$，有
\begin{equation}
P_{\pa_i,S}(v_i) = P_{\pa_i}(v_i). \tag{1.39}
\end{equation}

性质 1 使每个父节点集 $\PA_i$ 相对于其子节点 $V_i$ 成为外生的，从而确保条件概率 $P(v_i\mid\pa_i)$ 与“通过外部控制把 $\PA_i$ 置为 $\pa_i$”对 $V_i$ 所产生的效果相一致。性质 2 表达的是不变性：一旦控制了 $V_i$ 的直接原因 $\PA_i$，其他任何干预都不会影响 $V_i$ 的概率。

\subsection{因果关系及其稳定性}

这种以机制为基础的干预观念，为“因果效应”“因果影响”之类的概念提供了语义基础；这些概念将在第三、四章予以形式定义和分析。例如，要检验变量 $X_i$ 是否对另一变量 $X_j$ 有因果影响，我们用 \textup{(1.37)} 的截断分解公式计算 $X_j$ 在行动 $\dox{X_i=x_i}$ 之下的（边缘）分布，即对 $X_i$ 的一切取值 $x_i$ 计算 $P_{x_i}(x_j)$，并检验该分布是否敏感于 $x_i$。由前面的例子容易看出，只有 $X_i$ 在因果网络中的后代才可能受 $X_i$ 影响：从联合分布中删去因子 $P(x_i\mid\pa_i)$，会使 $X_i$ 在切除图中成为根节点，而根节点（正如 d-分离准则所规定的）与除其后代之外的一切变量独立。

对因果影响的这一理解，使我们能够确切地看出因果关系何以比概率关系更“稳定”，以及以何种方式更稳定。我们预期二者在稳定性上存在差异，是因为因果关系是本体论的，描述的是世界中客观的物理约束；而概率关系是认识论的，反映的是我们对世界的所知或所信。因此，只要环境未发生变化，即便我们关于环境的知识发生变化，因果关系也应保持不变。为说明这一点，考虑因果关系 $S_1$：“打开洒水器不会影响下雨”，并把它与其概率对应物 $S_2$ 相比：“洒水器的状态与雨的状态独立（或无关联）”。图 1.2 展示了 $S_2$ 发生变化而 $S_1$ 保持不变的两种明显方式。第一，当我们得知眼下是什么季节（$X_1$）时，$S_2$ 由假变真。第二，在已知季节的条件下，一旦观测到路面湿了（$X_4=\text{真}$），$S_2$ 又由真变假。另一方面，无论我们得知或知道了关于季节、关于路面的什么，$S_1$ 始终保持为真。

这个例子揭示了因果关系比相应概率关系更稳定的一种更强意义，它超出了二者在本体论—认识论上的基本差别。关系 $S_1$ 对于“季节如何影响洒水器”这一机制的变化保持不变的；事实上，它对于该因果图中所示一切机制的变化都保持不变。由此我们看到，因果关系对本体论变化也表现出更强的稳健性：它们所敏感的机制集合更小。更具体地说，与概率关系形成鲜明对照的是，因果关系对于支配因果变量的机制（本例中的 $X_3$）的变化保持不变。

鉴于这种稳定性，人们偏好以因果结构而非概率结构来编码知识，也就不足为奇了。概率关系，例如边缘独立与条件独立，或许有助于从不受控的观测中假设初始的因果结构；然而一旦知识被铸入因果结构，这些概率关系往往就被遗忘了：人们对某个领域中条件独立所作的种种判断，都是从已获得的因果结构派生出来的。这就解释了为什么人们能够自信地断言某些条件独立（例如中国的豆价与洛杉矶的交通独立），却对其中涉及的数值概率一无所知（例如豆价是否会超过每蒲式耳 10 美元）。

（机制的）稳定性要素也是所谓因果性的“解释论”说明的核心。依这种说明，因果模型不必编码干预之下的行为，而主要旨在提供对数据如何生成的“解释”或“理解”。\footnote{这种解释论说明的要素可见于 Dempster (1990)、Cox (1992) 与 Shafer (1996) 的著作；另见 King et al. (1994, p.~75)。}无论我们对事物的“理解”最终作何用途，我们确实更愿意以持久的、可跨情境迁移的关系来理解，而不是以转瞬即逝的关系为基础。伴随充分解释而来的那种“可理解感”，正是解释中所用因果关系之可迁移性（因而也是我们对它们的熟悉）的自然副产品。正是出于稳定性的考虑，我们认为气压计下降只能预测而不能解释下雨；这类预测无法迁移到气压计周围气压受人工控制的情境。真正的理解使人能够在此类新情境中作出预测——在那里，有些机制发生了变化，另一些机制被添加进来。因此，似乎有理由认为，在最终的分析中，因果性的解释论说明不过是操纵论说明的一个变体，只不过其中的干预是休眠的。据此，我们不妨把对“数据如何生成”“事物如何运作”的永不满足的追问，看作是在追求一种能力：在更广泛的情境中作出预测的能力，包括事物被拆解、被重构或经历自发变化的情境。

\section{函数因果模型}

我们引入贝叶斯网络之因果解释的方式，与因果模型（及因果图）最初被引入遗传学（Wright 1921）、计量经济学（Haavelmo 1943）和社会科学（Duncan 1975）的方式有根本不同，也与因果模型在物理学和工程学中的惯常用法不同。在那些模型中，因果关系以确定性函数方程的形式表达，概率则通过“方程中某些变量未被观测”这一假设引入。这反映了 Laplace (1814) 的自然现象观：自然规律是确定性的，随机性之所以出现，仅仅是因为我们对底层的边界条件无知。相反，因果贝叶斯网络定义中的一切关系都被假定为本质上随机的，因而诉诸现代（即量子力学的）物理观：自然的一切规律本质上都是概率性的，决定论只是一种方便的近似。

在本书中，我们将表示对 Laplace 准决定论因果观的偏好，并常用它（与随机论因果观相对照）来定义和分析我们所研究的大部分因果对象。这一偏好基于三点考虑。第一，Laplace 式的观念更具一般性。每个随机模型都可以由许多函数关系（带随机输入）来模拟，反之则不然；函数关系只能作为极限情形用随机模型来逼近。第二，Laplace 式的观念更契合人类直觉。少数与 Laplace 式观念之预测相冲突的玄奥量子力学实验，会引起惊讶与不信，并要求物理学家放弃关于定域性与因果性的根深蒂固的直觉（Maudlin 1994）。我们的目标是保存、阐明并满足——而不是摧毁——这些直觉。\footnote{常听到的一种论调是：人类直觉属于心理学，而不属于科学或哲学。但在因果直觉的问题上，这种论调并不适用——因果思想的创始者们本就不能被忽视，概念的含义受到质疑时更是如此。事实上，是否合乎人类直觉，一直是每一种因果哲学研究之充分性的终极判据；而要把背景信息恰当地纳入统计研究，同样有赖于对因果判断的准确解释。}

第三，有些在人类话语中无处不在的概念，只能在 Laplace 式框架中定义。例如我们将会看到，“事件 $B$ 之所以发生是因为事件 $A$”的概率，以及“若不是事件 $A$，事件 $B$ 本会不同”的概率，这类简单概念无法用纯随机模型来定义。这些所谓的反事实概念，需要把 Laplace 模型中体现的确定性成分与概率成分综合起来。

\subsection{结构方程}

函数因果模型的一般形式由一组如下形状的方程构成：
\begin{equation}
x_i = f_i(\pa_i, u_i),\qquad i=1,\dots,n, \tag{1.40}
\end{equation}
其中 $\pa_i$（意为父节点）表示直接决定 $X_i$ 之取值的变量集，$U_i$ 表示因省略因素而产生的误差（或“扰动”）。式 \textup{(1.40)} 是线性结构方程模型（SEM）
\begin{equation}
x_i = \sum_{k\ne i} \beta_{ik}\,x_k + u_i,\qquad i=1,\dots,n, \tag{1.41}
\end{equation}
的非线性、非参数推广；后者已成为经济学与社会科学中的标准工具（关于这一事业的详细阐述见第五章）。在线性模型中，$\pa_i$ 对应于 \textup{(1.41)} 右端系数非零的那些变量。

\textup{(1.40)} 中函数关系的解释，就是函数在物理学与自然科学中所具有的标准解释：它是一条配方、一种策略或一条定律，规定了面对 $(\PA_i,U_i)$ 可能取得的每一种取值组合，自然将给 $X_i$ 赋予什么值。一组形如 \textup{(1.40)} 的方程，若每个方程都代表一个自主的机制，则称为\emph{结构模型}；若每个变量都有一个以它作为左端（称为因变量）的独立方程，则称该模型为\emph{结构因果模型}，或简称\emph{因果模型}。\footnote{关于因果模型、结构方程与误差项的形式处理，见第五章（\S 5.4.1）与第七章（\S 7.1、\S 7.2.5）。}在数学上，结构方程与代数方程的区别在于：结构方程的任何子集本身就是一个有效的结构模型——它表示某一组干预之下的情形。

为作说明，图 1.5 描绘了一个典型的计量经济学模型，它通过下列方程把价格与需求联系起来：
\begin{equation}
q = b_1 p + d_1 i + u_1, \tag{1.42}
\end{equation}
\begin{equation}
p = b_2 q + d_2 w + u_2, \tag{1.43}
\end{equation}
其中 $Q$ 是家庭对产品 $A$ 的需求量，$P$ 是产品 $A$ 的单价，$I$ 是家庭收入，$W$ 是生产产品 $A$ 的工资率，$U_1$ 与 $U_2$ 是误差项——分别表示影响数量与价格的未建模因素（Goldberger 1992）。与该模型相联系的图是循环的，与变量 $U_1,U_2,I,W$ 相联系的顶点都是根节点，传达出它们相互独立的假设。自主性（Aldrich 1989）在这里意味着：两个方程代表经济中两个松散耦合的部分——消费者与生产者。式 \textup{(1.42)} 描述消费者如何决定购买多少数量 $Q$，\textup{(1.43)} 描述厂商如何决定索取什么价格 $P$。如同一切反馈系统，这里也隐含着动态：今天的价格依据昨天的需求确定，而这些价格又将决定下一期交易的需求。在背景量 $U_1,U_2$ 保持不变的假设下，此类方程的解表示一种长期均衡。

相对于变化的动态，这两个方程被认为是“自主的”，意思是影响其中一个方程的外部变化并不意味着其他方程也要变化。例如，若政府决定实行价格管制，把价格 $P$ 定为 $p_0$，则 \textup{(1.43)} 将修改为 $p=p_0$，而 \textup{(1.42)} 中的关系保持不变，于是 $q=b_1p_0+d_1i+u_1$。由此我们看到，$b_1$ 这一“需求弹性”应被解释为 $P$ 每单位\emph{受控}变化所引起的 $Q$ 的变化率。这当然不同于 $P$ 每单位\emph{观测}变化（在不受控条件下）所引起的 $Q$ 的变化率；后者除 $b_1$ 之外，还受 \textup{(1.43)} 各参数的影响（见 \S 7.2.1 式 \textup{(7.14)}）。受控变化与观测变化之别，对于正确解释社会科学与经济学中的结构方程模型至关重要，第五章将详加讨论。若有理由相信在价格管制政策下消费者行为也会改变，那么这种改变后的行为需要显式建模——例如，把系数 $b_1$ 与 $d_1$ 当作涉及 $P$ 的辅助方程中的因变量。\footnote{确实，消费者通常会对价格管制作出反应：预期短缺而囤积商品（Lucas 1976）。不过这类现象对结构模型并不陌生，只是需要用更精细的方程来刻画消费者的预期。}\S 7.2.1 将用该模型分析政策相关问题。

为说明非线性函数模型的运作，再次考虑图 1.2 所描绘的因果关系。与这些关系相联系的因果模型将由五个函数构成，每个函数代表支配一个变量的自主机制：
\begin{equation}
x_1=u_1,\quad x_2=f_2(x_1,u_2),\quad x_3=f_3(x_1,u_3),\quad x_4=f_4(x_3,x_2,u_4),\quad x_5=f_5(x_4,u_5). \tag{1.44}
\end{equation}
误差变量 $U_1,\dots,U_5$ 未在图中显式画出；按惯例，这意味着假定它们相互独立。当某些扰动被判断为相依时，习惯上通过在图中添加双箭头来编码这种相依，如图 1.1(a) 所示。

函数 $\{f_1,\dots,f_5\}$ 与扰动项的一个典型指定由下述布尔模型给出：
\begin{equation}
\begin{aligned}
x_2 &= [(X_1=\text{冬})\vee(X_1=\text{秋})\vee u_2]\wedge\neg u_2',\\
x_3 &= [(X_1=\text{夏})\vee(X_1=\text{春})\vee u_3]\wedge\neg u_3',\\
x_4 &= (x_2\vee x_3\vee u_4)\wedge\neg u_4',\\
x_5 &= (x_4\vee u_5)\wedge\neg u_5',
\end{aligned} \tag{1.45}
\end{equation}
其中 $x_i$ 表示 $X_i=\text{真}$，$u_i$ 与 $u_i'$ 分别表示触发性的与抑制性的异常。例如，$u_4$ 表示在洒水器关闭（$\neg x_3$）且不下雨（$\neg x_2$）时仍可能使路面变湿（$X_4$）的（未指明的）事件（如水管破裂）；而 $u_4'$ 则表示尽管下雨（$x_2$）、洒水器开着（$x_3$）且有 $u_4$，仍会使路面保持干燥的（未指明的）事件（如路面盖着塑料布）。

必须强调，在刚才描述的两个模型中，位于等号左端的变量（因变量或输出变量）在各自方程中的作用与其他变量不同。这一区分的意义在讨论干预时就清楚了：正是通过这一区分，我们才能确定在“把价格定为 $p_0$”（$\dox{P=p_0}$）或“打开洒水器”（$\dox{X_3=\text{真}}$）这类局部干预之下，应当修改哪个方程。\footnote{有些经济学家把供求方程写成 $\{q=ap+u_1,\ q=bp+u_2\}$，让 $q$ 同时出现在两个方程的右端；这样做就放弃了分析价格管制政策的可能，除非另有一套符号机制来识别 $\doxp(P=p_0)$ 算子将修改哪个方程。}

下面我们把 \textup{(1.40)} 所定义的函数模型的特征，与 \S 1.3 所定义的因果贝叶斯网络的特征作一比较。为此，我们考虑三类查询的处理：
\begin{itemize}
\item \emph{预测}（例如：若发现洒水器关着，路面会打滑吗？）；
\item \emph{干预}（例如：若我们确保洒水器关着，路面会打滑吗？）；
\item \emph{反事实}（例如：已知路面实际上没有打滑而洒水器开着，那么若当时洒水器关着，路面会打滑吗？）。
\end{itemize}
我们将会看到，这三类查询构成一个由三种根本不同的问题组成的层级，对知识细节的要求逐级提高。

\begin{figure}[htbp]
\centering
\begin{tikzpicture}[causal,scale=1.0,every node/.style={transform shape}]
  \node[vtx] (U1) at (-3.4,2.2) {$U_1$};
  \node[vtx] (I)  at (-2.2,2.2) {$I$};
  \node[vtx] (W)  at (2.2,2.2) {$W$};
  \node[vtx] (U2) at (3.4,2.2) {$U_2$};
  \node[vtx] (Q)  at (-1.6,0.5) {$Q$};
  \node[vtx] (P)  at (1.6,0.5) {$P$};
  \draw[ed] (U1) -- (Q);
  \draw[ed] (I)  -- node[above,font=\scriptsize,pos=0.55]{$d_1$} (Q);
  \draw[ed] (W)  -- node[above,font=\scriptsize,pos=0.55]{$d_2$} (P);
  \draw[ed] (U2) -- (P);
  \draw[ed] (P) to[bend left=28] node[above,font=\scriptsize,pos=0.5]{$b_2$} (Q);
  \draw[ed] (Q) to[bend left=28] node[below,font=\scriptsize,pos=0.5]{$b_1$} (P);
\end{tikzpicture}
\caption{说明价格（$P$）、需求（$Q$）、收入（$I$）与工资（$W$）之间关系的因果图。}
\label{fig:1.5}
\end{figure}

\subsection{因果模型中的概率预测}

给定一个因果模型（式 \textup{(1.40)}），若从 $\PA_i$ 的每个成员向 $X_i$ 画箭头，所得的图 $G$ 称为\emph{因果图}。若因果图无环，则相应的模型称为\emph{半马尔可夫的}，此时 $X$ 变量的值由 $U$ 变量的值唯一确定。在此条件下，联合分布 $P(x_1,\dots,x_n)$ 由误差变量的分布 $P(u)$ 唯一确定。若除无环之外，误差项还联合独立，则称该模型为\emph{马尔可夫的}。

关于马尔可夫模型的一个基本定理，通过定理 1.2.7 的父节点马尔可夫条件，建立了因果性与概率之间的联系。

\begin{theorem}[因果马尔可夫条件]
每个马尔可夫因果模型 $M$ 都诱导出一个分布 $P(x_1,\dots,x_n)$，它相对于与 $M$ 相联系的因果图 $G$ 满足父节点马尔可夫条件；也就是说，每个变量 $X_i$ 在给定其在 $G$ 中的父节点 $\PA_i$ 时，与它的所有非后代独立（Pearl and Verma 1991）。\footnote{考虑到该定理的一般性与透明性，若某一版本早已见于文献，我不会感到意外；但我并不知道有任何非参数版本。}
\end{theorem}

证明是直接的。既然集合 $\{\PA_i,U_i\}$ 确定了 $X_i$ 的一个唯一取值，分布 $P(x_1,\dots,x_n,u_1,\dots,u_n)$ 当然相对于增广 DAG $G(X,U)$ 是马尔可夫的，其中 $U$ 变量被显式表示。所需的边缘分布 $P(x_1,\dots,x_n)$ 的马尔可夫条件，则由 $G(X,U)$ 中的 d-分离得出。

定理 1.4.1 表明，定理 1.2.7 的父节点马尔可夫条件可由两条因果假设推出：(1) 我们承诺把凡是两个或更多其他变量之原因的变量都纳入模型（而非留作背景）；(2) Reichenbach (1956) 的共同原因假设，也称“无因果则无相关”：若两个变量相依，则其中之一是另一个的原因，或者存在第三个变量同时是二者的原因。这两条假设意味着 $U$ 中的背景因素相互独立，因而因果模型是马尔可夫的。定理 1.4.1 既解释了为什么马尔可夫模型在因果分析中如此常被假定，也解释了为什么父节点马尔可夫条件（定理 1.2.7）常被视为因果模型的固有特征（如见 Kiiveri et al. 1984; Spirtes et al. 1993）。\footnote{Kiiveri 等 (1984) 那篇题为“递归因果模型”的论文，首次（对严格正分布）证明了定理 1.2.7 的父节点马尔可夫条件可由 \textup{(1.33)} 的分解推出。然而这一蕴含纯粹是概率性的，并未援引因果性的任何方面。要在因果性与概率之间建立联系，我们必须先为因果性设计一个模型，或者用操纵（如定义 1.3.1），或者用结构方程中的函数关系（如定理 1.4.1）。}

因果马尔可夫条件意味着：把每个子—父关系刻画为确定性函数，而不是通常的条件概率 $P(x_i\mid\pa_i)$，对所得分布施加的独立约束是等价的，并导致与贝叶斯网络相同的递归分解（见式 \textup{(1.33)}）。更重要的是，这一点与函数 $\{f_i\}$ 的选择无关，也与误差分布 $P(u_i)$ 无关。因此，我们不必事先指定 $\{f_i\}$ 的函数形式或 $P(u_i)$ 的分布；一旦测量（或估计）了 $P(x_i\mid\pa_i)$，马尔可夫因果模型的一切概率性质便都确定了，而不管究竟是哪个机制在生成这些条件概率。Druzdzel and Simon (1993) 证明，对每个以分布 $P$ 为特征的贝叶斯网络 $G$（如 \textup{(1.33)}），都存在一个函数模型（如 \textup{(1.40)}）生成与 $P$ 完全相同的分布。\footnote{第九章将证明，除某些病态情形外，具有该性质的函数模型实际上有无穷多个。}由此可知，在贝叶斯网络的一切概率应用中——包括统计估计、预测与诊断——我们都可以使用 \textup{(1.40)} 所指定的等价函数模型，并可以把函数模型仅仅看作编码联合分布函数的另一种方式。

尽管如此，即便在纯预测（即非操纵）任务中，因果—函数式的指定也比概率式的指定有若干优点。首先也最重要的是，因果图 $G$ 所示的一切条件独立关系都保证是稳定的——即对函数 $f_i$ 与分布 $P(u_i)$ 所代表机制的参数变化保持不变。这意味着，选择用马尔可夫因果模型组织知识的智能体，能够在不评估数值概率的情况下对条件独立关系作出可靠断言——这是人类常有的能力，也是推断的一项有用特性。第二，函数式指定往往更有意义、更自然，并且只产生少量参数。典型例子是社会科学与经济学中使用的线性结构方程（见第五章），以及在刻画多个二值原因之效应时颇为流行的“噪声或门”（Pearl 1988b, p.~184）。第三（也许最难为经验主义者接受），在函数模型中，关于可观测量之间条件独立的判断性假设被简化、也变得更可靠，因为这类假设直接被表述为关于未观测共同原因存在与否的判断（例如：为什么中国的豆价会被判断为与洛杉矶的交通独立？）。例如，在构造贝叶斯网络时，我们不必逐一判断每个变量是否（在给定父节点时）与其所有非后代独立，而只需判断父节点集是否包含了全部相关的直接原因——特别是，判断父节点集中被遗漏的因素是否构成另一个被观测变量的原因。这类判断更为自然，因为它们可以直接从定性的因果结构中识别出来，而这一结构正是我们心智为储存经验的稳定方面而选定的。

最后，把预测模型建立在因果机制之上还有一个优点，它来自关于稳定性的考虑（\S 1.3.2）。当环境中的某些条件发生变化时，通常只有少数因果机制受其影响，其余保持不变。此时，明知相应的符号改动也是局部的、只涉及少数参数，去（以判断方式）重新评估或（以统计方式）重新估计模型参数，总比从头重新估计整个模型简单。\footnote{据我所知，因果模型的这一方面尚未被形式地研究过；此处提出，作为适应性系统研究者的一个研究课题。}

\subsection{函数模型中的干预与因果效应}

函数式刻画 $x_i=f_i(\pa_i,u_i)$ 与它的随机对应物一样，提供了一种便利的语言，用以指定所得分布将如何响应外部干预而变化。其做法是把每个干预编码为一组精选函数的改动，而不是一组精选条件概率的改动。于是，干预的总体效果便可通过修改模型中相应的方程、并用修改后的模型计算新的概率函数来预测。因此，因果贝叶斯网络（\S 1.3）的一切特征都可以在马尔可夫函数模型中模拟。

例如，要在 \textup{(1.44)} 的模型中表示“打开洒水器”这一行动，我们删去方程 $x_3=f_3(x_1,u_3)$，代之以 $x_3=\mathrm{On}$。修改后的模型包含计算该行动对其他变量之一切效果所需的全部信息。例如，修改后的模型所诱导的概率函数将等于 \textup{(1.36)} 所给者，修改后的图也将与图 1.4 一致。

更一般地，当干预迫使变量子集 $X$ 取定值 $x$ 时，就要从 \textup{(1.40)} 的模型中剪除一个方程子集——$X$ 的每个成员对应一个——从而在其余变量上定义一个新的分布，它刻画了干预的效果，并与从因果贝叶斯网络中剪除家族所得的截断分解（式 \textup{(1.37)}）一致。\footnote{把干预显式翻译为从模型中“抹去”方程的做法，最早由 Strotz and Wold (1960) 提出，后为 Fisher (1970) 与 Sobel (1990) 所用。更为精细的干预类型，涉及条件行动与随机策略，将在第四章中表述。}

函数模型对干预的表示比随机模型更灵活、更一般。第一，干预分析可以推广到循环模型，如图 1.5 中的模型，从而回答与政策有关的问题\footnote{这类问题，尤其是涉及内生变量控制的问题，在计量经济学教科书中显著缺席（见第五章）。}（例如：若我们把价格控制在 $p_0$，需求量将是多少？）。第二，涉及方程参数改动（如 \textup{(1.42)} 中的 $b_1$ 与 $d_1$）的干预，比那些被描述为条件概率修改者的干预更容易理解，这也许是因为稳定的物理机制通常与方程相联系，而不与条件概率相联系。条件概率被认为是可从联合分布导出的，而不是联合分布的生成者。第三，非马尔可夫模型中因果效应的分析，将因使用函数模型而大大简化。原因在于：条件概率 $P(x_i\mid\pa_i)$ 有无穷多个，而在离散变量 $X_i$ 与 $\PA_i$ 之间，函数 $x_i=f_i(\pa_i,u_i)$ 却只有有限多个。这一事实将使我们在第八章（\S 8.2.2）中能够用线性规划技术，对涉及不依从（noncompliance）的研究中的因果效应求得锐界。

最后，函数模型允许分析情境特定的行动与政策。迄今为止所定义的因果效应概念，在实践的政策制定中用途有限。原因在于，因果效应告诉我们某一行动引起反应的一般倾向（例如某药物在总体人群中促进康复的倾向），但并不针对在由一组特定观测所刻画、而这些观测本身又可能受该行动影响的情境中的行动。医生通常关心的是治疗对某位已经检查、发现具有某些症状的病人的效果；其中有些症状本身会受到治疗的影响。同样，经济学家关心的是在由各种经济指标刻画的特定经济情境中征税的效果，而这些指标（同样）在征税实施后也会受到影响。这类情境特定的因果效应，无法通过在静态贝叶斯网络中模拟干预来计算，因为情境本身随干预而变，从而条件概率 $P(x_i\mid\pa_i)$ 在这个过程中被改变了。然而，函数关系 $x_i=f_i(\pa_i,u_i)$ 保持不变，这使我们能够计算情境特定的因果效应，下一节将作概述（详见 \S 7.2.1、\S 8.3 与 \S 9.3.4）。

\subsection{函数模型中的反事实}

现在转向函数模型最独特的特征——反事实分析。如前所述，某些反事实语句无法在随机因果网络的框架中定义。为看清困难所在，让我们考虑最简单的因果贝叶斯网络：一对相互独立（因而互不连通）的二值变量 $X$ 与 $Y$。例如，在对照（即随机化）临床试验中，若发现治疗 $X$ 对受试者反应 $Y$ 的分布没有影响，就会出现这样的网络；$Y$ 可以表示康复（$Y=0$）或死亡（$Y=1$）。假定某位受试者乔接受了治疗并死亡；我们要问：乔的死是因为治疗、尽管有治疗，还是与治疗无关？换言之，我们要求的是“若乔未接受治疗，他会死亡”的概率 $Q$。

为凸显回答此类反事实问题的困难，让我们取一个极端的例子：在治疗组与对照组中，均有 50\% 的病人康复、50\% 死亡；并进一步假定样本量趋于无穷，于是
\begin{equation}
P(y\mid x) = 1/2 \qquad\text{对一切 } x \text{ 与 } y. \tag{1.46}
\end{equation}
熟悉统计检验的读者立刻会看出，仅凭现有数据无法回答这个反事实问题：乔接受了治疗并死亡，从未在无治疗条件下被检验过。此外，困难并不源于把问题指向某个特定个人乔——对他我们只有一个数据点。若把问题改述为总体频率——在治疗下死亡的受试者中，有多大比例 $Q$ 若未接受治疗本会康复——仍会遇到同样的困难，因为这些人中没有谁在无治疗条件下被检验过。这类困难促使一些统计学家把反事实问题斥为形而上学，并主张把统计分析限制在能由直接检验回答的问题上（Dawid 2000）。

然而，我们的科学语言、法律语言与日常语言中充斥着反事实话语，这清楚表明反事实远非形而上学；它们必定具有确定的可检验含义，并承载着有价值的实质信息。因此，反事实分析对任何认同本书宗旨的人而言都是一个机会：把实质知识与统计数据结合起来，以精炼前者、解释后者。在这一框架内，反事实问题要求回答一些棘手但可以处理的技术问题：反事实查询的经验内容是什么？回答这些查询需要什么知识？这些知识如何用数学表示？给定这样的表示，又需要什么数学机械来导出答案？

第七章（\S 7.2.2）把反事实阐发为关于某些机制在时间上持存的断言。在我们的例子中，假定每位（存活）病人对治疗的反应是持存的。若结局 $Y$ 是一种可逆状态而非死亡，那么反事实断言将直接转化为关于未来治疗反应的预测。但即便在死亡的情形，反事实量 $Q$ 也不仅仅意味着对已死亡受试者假想行为的猜测，它还包含一个关于未治疗的存活受试者在后续治疗下如何的可检验断言。我们把它留作读者的练习：证明在 \textup{(1.46)} 的基础上、排除抽样波动，治疗组中死亡的受试者里“若未接受治疗本会康复”的比例 $Q$，恰好等于未治疗组中存活的受试者里“若给予治疗将会死亡”的比例 $Q'$。\footnote{例如，若 $Q$ 等于 100\%（即所有接受治疗并死亡者若未接受治疗本会康复），则未治疗组中所有存活者若给予治疗都将死亡（同样排除抽样波动）。当我们建立起分析因果概率的数学机械之后，这类练习将变得寻常（见第九章定理 9.2.12、式 \textup{(9.11)}—\textup{(9.12)}）。} $Q$ 是假设性的，$Q'$ 则毫无疑问是可检验的。

在勾勒出反事实的经验解释之后，本章下一步要问的是表示问题：回答反事实问题需要什么知识？这些知识应如何表述，才能使反事实查询得以迅速而可靠地回答？这种表示确实存在，这一点从人们区分可信与不可信反事实语句的迅捷与一致中即可看出。多数人会同意：若克林顿未曾遇见莫妮卡·莱温斯基，他在历史上的地位会有所不同；但只有少数人会断言：若他昨天没吃早饭，他的历史地位就会改变。在认知科学中，这种意见的一致性已近乎一种证明：某种有效表示和操作反事实的机械，就存在于人类心智的某处。那么，这一机械的构件是什么呢？

一种直截了当的表示方案是：(i) 把反事实知识以反事实前提的形式储存起来；(ii) 用某些能够从前提推出结论的逻辑推理规则，导出反事实查询的答案。哲学家 Robert Stalnaker (1968) 与 David Lewis (1973a,b) 确实采取了这种做法，他们用最接近世界语义学构造了反事实逻辑（即：“若 $A$ 则 $B$”为真，当且仅当在（距我们）最接近的、$A$ 为真的可能世界中 $B$ 为真）。然而，最接近世界语义学仍留下两个未答的问题。(1) 选择什么样的距离度量，才能使反事实推理与通常的因果观念相容？(2) 什么样的世界间距离的心智表示，才能使反事实的计算（对人与机器而言）变得可操作、可实践？这两个问题由第七章展开的结构模型进路予以回答。

统计学家在潜在结局框架中采取了一条与 Lewis 类似（尽管不那么形式化）的进路（Rubin 1974; Robins 1986; Holland 1988）。在这里，实质知识以反事实变量之间概率关系（如独立）的形式表达，然后用于估计因果效应。表示问题便从最接近世界进路转移到潜在结局进路：反事实之间的概率关系在研究者心中是如何储存或推断的？在第七章（另见 \S 3.6.3）中，我们将分析最接近世界进路与潜在结局进路，并把它们与结构模型进路相比较；在后者中，反事实由函数因果模型（式 \textup{(1.40)}）导出（事实上由之定义）。

为看清反事实与结构方程之间的联系，我们应先考察：为什么贝叶斯网络所编码的信息，即便在其因果解释下，也不足以回答反事实查询。再次考虑对照随机实验的例子（式 \textup{(1.46)}），它对应于一个无边的贝叶斯网络（图 1.6(a)），含两个相互独立的二值变量，其联合概率为
\begin{equation}
P(y,x) = 0.25 \qquad\text{对一切 } x \text{ 与 } y. \tag{1.47}
\end{equation}
下面给出两个函数模型，它们都生成 \textup{(1.47)} 的联合概率，却给所关心的量 $Q$ 赋予不同的值——$Q$ 即“在治疗下死亡（$x=1,y=1$）的受试者，若未接受治疗（$x=0$）本会康复（$y=0$）”的概率。

\medskip
\noindent\emph{模型 1}（图 1.6(b)）\quad 令
\[
x = u_1,\qquad y = u_2,
\]
其中 $U_1$ 与 $U_2$ 是两个相互独立的二值变量，$P(u_1=1)=P(u_2=1)=\tfrac{1}{2}$（例如随机掷硬币）。

\medskip
\noindent\emph{模型 2}（图 1.6(c)）\quad 令
\begin{equation}
x = u_1,\qquad y = x\,u_2 + (1-x)(1-u_2), \tag{1.48}
\end{equation}
其中 $U_1$ 与 $U_2$ 同前，仍是两个相互独立的二值变量。

模型 1 对应于这样一种治疗（$X$）：它对任何受试者都没有效果；在模型 2 中，每个受试者都受治疗影响。两个模型产生相同分布的原因在于，模型 2 描述的是两个子总体的混合。在一个子总体（$u_2=1$）中，每个受试者当且仅当接受治疗时死亡（$y=1$）；在另一个子总体（$u_2=0$）中，每个受试者当且仅当接受治疗时康复（$y=0$）。对应于这两个模型的分布 $P(x,y,u_2)$ 与 $P(x,y)$ 见图 1.7。

$Q$ 的值在这两个模型中不同。在模型 1 中，$Q$ 求值为零，因为死亡的受试者对应于 $u_2=1$，而由于治疗对 $y$ 没有影响，把 $X$ 从 1 改为 0 仍将得到 $y=1$。然而在模型 2 中，$Q$ 求值为 1，因为在治疗下死亡的受试者必定对应于 $u_2=1$（即那些一治就死的人），这意味着他们当且仅当未接受治疗时才会康复。

\begin{figure}[htbp]
\centering
\begin{tikzpicture}[causal,scale=1.0,every node/.style={transform shape}]
  % (a)
  \begin{scope}
    \node[circle,fill,inner sep=1.4pt] (Xa) at (0,0.8) {};
    \node[font=\small,above=1.2mm] at (0,0.8) {$X$};
    \node[circle,fill,inner sep=1.4pt] (Ya) at (0.9,-0.6) {};
    \node[font=\small,below=1.2mm] at (0.9,-0.6) {$Y$};
    \node[font=\small] at (0.45,-1.7) {(a)};
  \end{scope}
  % (b)
  \begin{scope}[xshift=3.6cm]
    \node[vtx] (Xb) at (0,0) {$X$};
    \node[vtx] (Yb) at (2.0,0) {$Y$};
    \node[lat] (U1b) at (-0.3,1.35) {$U_1$};
    \node[lat] (U2b) at (2.3,1.35) {$U_2$};
    \draw[ed] (Xb) -- (Yb);
    \draw[edd] (U1b) -- (Xb);
    \draw[edd] (U2b) -- (Yb);
    \node[font=\small] at (1.0,-1.1) {(b)};
  \end{scope}
  % (c)
  \begin{scope}[xshift=8.6cm]
    \node[vtx] (Xc) at (0,0) {$X$};
    \node[vtx] (Yc) at (2.0,0) {$Y$};
    \node[lat] (U1c) at (-0.3,1.35) {$U_1$};
    \node[lat] (U2c) at (2.3,1.35) {$U_2$};
    \draw[edb] (Xc) -- (Yc);
    \draw[edd] (U1c) -- (Xc);
    \draw[edd] (U2c) -- (Yc);
    \node[font=\small] at (1.0,-1.1) {(c)};
  \end{scope}
\end{tikzpicture}
\caption{(a) 表示 \textup{(1.47)} 之分布的因果贝叶斯网络。(b) 按模型 1 表示生成 (a) 中分布之过程的因果图。(c) 按模型 2 表示的同一过程。（$U_1$ 与 $U_2$ 均未被观测。）}
\label{fig:1.6}
\end{figure}

\begin{figure}[htbp]
\centering
\small
\begin{tabular}{@{}lcccccc@{}}
\toprule
& \multicolumn{2}{c}{$u_2=0$} & \multicolumn{2}{c}{$u_2=1$} & \multicolumn{2}{c}{边缘} \\
\cmidrule(lr){2-3}\cmidrule(lr){4-5}\cmidrule(l){6-7}
& $x=1$ & $x=0$ & $x=1$ & $x=0$ & $x=1$ & $x=0$ \\
\midrule
\multicolumn{7}{@{}l}{\textit{模型 1}}\\
$y=1$（死亡）   & 0    & 0    & 0.25 & 0.25 & 0.25 & 0.25\\
$y=0$（康复）   & 0.25 & 0.25 & 0    & 0    & 0.25 & 0.25\\[2pt]
\multicolumn{7}{@{}l}{\textit{模型 2}}\\
$y=1$（死亡）   & 0    & 0.25 & 0.25 & 0    & 0.25 & 0.25\\
$y=0$（康复）   & 0.25 & 0    & 0    & 0.25 & 0.25 & 0.25\\
\bottomrule
\end{tabular}
\caption{列联表：两个模型的分布 $P(x,y,u_2)$ 与 $P(x,y)$。}
\label{fig:1.7}
\end{figure}

这个例子的第一点教益是：随机因果模型不足以计算反事实的概率；为进行计算，需要关于 $P(y\mid x)$ 背后实际过程的知识。\footnote{在潜在结局框架中（\S 3.6.3 与 \S 7.4.4），这类知识通过定义反事实变量 $Y_1$ 与 $Y_0$ 上的分布而获得随机的外观；$Y_1$、$Y_0$ 分别表示个体对治疗与无治疗的潜在反应。这些假设性变量所起的作用类似于我们模型中的函数 $f_i(\pa_i,u_i)$；它们体现了这样一个确定性假设：每个个体对治疗都有确定的反应，无论该治疗是否真的实现。}第二点教益是：函数因果模型构成了足以计算（并定义）此类概率的数学对象。例如，考虑 \textup{(1.48)} 的模型 2。我们据以断定“死亡的受治者（$y=1,x=1$）若未受治本会康复”的方式，包含三个心智步骤。第一，把手头的证据 $e:\{y=1,x=1\}$ 应用于模型，断定 $e$ 只与 $U_1$、$U_2$ 的一种实现相容，即 $\{u_1=1,u_2=1\}$。第二，为模拟“若他或她未接受治疗”这一假设条件，我们把 $x=0$ 代入 \textup{(1.48)}，同时忽略第一个方程 $x=u_1$。最后，在 $x=0$、$u_2=1$ 的假定下解 \textup{(1.48)} 求 $y$，得到 $y=0$，据此断定在所述假设条件下康复（$y=0$）的概率为 1。

这三个步骤可以推广到任一因果模型 $M$，如下所述。给定证据 $e$，为计算在假设条件 $X=x$（$X$ 为变量子集）下 $Y=y$ 的概率，对 $M$ 施行以下三步。

\begin{description}[leftmargin=0pt,style=nextline,itemsep=0.2em]
\item[\textbf{步骤 1（回溯）}] 更新概率 $P(u)$，得到 $P(u\mid e)$。
\item[\textbf{步骤 2（行动）}] 把集合 $X$ 中变量所对应的方程替换为方程 $X=x$。
\item[\textbf{步骤 3（预测）}] 用修改后的模型计算 $Y=y$ 的概率。
\end{description}

用时间的隐喻来说，这三步程序可解释如下。步骤 1 依据当前证据 $e$ 解释过去（$U$）；步骤 2（最小程度地）扭转历史进程，使之符合假设条件 $X=x$；最后，步骤 3 依据我们对过去的新理解与新建的条件 $X=x$ 预测未来（$Y$）。

回想对 $U$ 的每个取值 $u$，$Y$ 都有唯一解，则步骤 3 总能给出所需概率的唯一解便显而易见：只需把所有使 $Y=y$ 成为解的 $u$ 所对应的概率 $P(u\mid e)$ 相加。第七章将发展计算反事实概率的有效程序，这些程序基于“孪生”网络中的概率传播（Balke and Pearl 1995a）：一个网络表示实际世界，另一个表示反事实世界。

注意，假设条件 $X=x$ 总是与所考虑模型中 $U$ 的现行取值 $u$ 相矛盾（否则 $X=x$ 就会真的实现，因而不会被当作假设）。正因如此，我们才在步骤 2 中援引外部干预（或者叫“理论改变”或“奇迹”，Lewis 1973b），它修改模型，从而把矛盾解释掉。第七章将扩展这一结构—干预模型，为反事实及反事实的概率提供完整的语义与公理化说明。与 Lewis 的理论不同，这一说明不基于假设世界之间相似性的抽象概念，而是立足于产生所考虑假设世界的实际机制。同样，与潜在结局框架不同，结构论说明不把反事实当作未定义的原始概念，而当作可从因果机制及其结构这些更基本的概念中导出的量。

反事实推理的三步模型还揭示了随机因果模型何以不足以计算反事实概率的真正原因。由于 $U$ 变量不在随机模型中显式出现，我们无法施行步骤 1，用手头证据 $e$ 更新 $P(u)$。这意味着若干基于反事实的常见概念——包括因果概率（已知结果）、解释的概率、以及情境相依的因果效应——都无法在这样的模型中定义。为此，我们必须对函数 $f_i$ 的形式与误差项的概率作出某些假设。例如，线性性、正态性与误差独立这三条假设，足以计算图 1.5 模型中的一切反事实查询（见 \S 7.2.1）。第九章将给出一些条件，在这些条件下，当 $f_i$ 与 $P(u)$ 未知、只假定其一般特征（如单调性）时，关于因果概率的反事实查询仍可由数据推断。同样，第八章（\S 8.3）将给出在只有随机模型可用时对反事实概率作界的方法。

上述考虑进一步表明，本节开头所列的三项任务——预测、干预与反事实——构成因果推理任务的一个自然层级，其精细化程度逐级提高，对完成任务所需知识的要求也逐级提高。预测是三者中最简单的，只需指定一个联合分布函数。干预分析除联合分布外，还需要因果结构。最后，处理反事实是最难的任务，因为它需要关于函数关系和（或）被省略因素之分布的某些信息。

这一层级也划定了本书各章的自然分工。第二章主要处理因果贝叶斯网络的概率方面（尽管底层的因果结构将充当概念指引）。第三至六章专门处理因果模型的干预方面，包括因果效应的识别、结构方程模型的澄清，以及混杂与可压缩性之间的关系。第七至十章处理反事实分析，包括公理化基础、在政策分析中的应用、反事实查询的定界、因果概率的识别，以及单一事件因果性的阐发。

我祝愿读者在这几章中经历一段顺畅而有益的旅程。但首先，还有一个重要的停靠站：术语的区分。
